\documentclass{article}

\usepackage{amsmath, amsthm, amssymb, amsfonts}
\usepackage{thmtools}
\usepackage{graphicx}
\usepackage{setspace}
\usepackage{geometry}
\usepackage{float}
\usepackage[backend=biber,style=ieee,sorting=ynt]{biblatex}
\usepackage[hidelinks]{hyperref}
\usepackage[utf8]{inputenc}
\usepackage[spanish]{babel}
\usepackage{framed}
\usepackage[dvipsnames]{xcolor}
\usepackage{tcolorbox}
\usepackage{tikz}
\usetikzlibrary{positioning}
\usetikzlibrary{arrows.meta}
\usepackage{caption}
\usepackage{longtable}
\usepackage{pdflscape}
\usepackage{svg}
\usepackage{subcaption}
\usepackage{caption}
\usepackage{multirow}
\usepackage{array}
\usepackage{listings}
\usepackage{cancel}
\usepackage{fancyhdr}

\addbibresource{referencias.bib}


\colorlet{LightGray}{White!90!Periwinkle}
\colorlet{LightOrange}{Orange!15}
\colorlet{LightGreen}{Green!15}



\newcommand{\HRule}[1]{\rule{\linewidth}{#1}}

\declaretheoremstyle[name=Theorem,]{thmsty}
\declaretheorem[style=thmsty,numberwithin=section]{theorem}
\tcolorboxenvironment{theorem}{colback=LightGray}

\declaretheoremstyle[name=Proposition,]{prosty}
\declaretheorem[style=prosty,numberlike=theorem]{proposition}
\tcolorboxenvironment{proposition}{colback=LightOrange}

\declaretheoremstyle[name=Principle,]{prcpsty}
\declaretheorem[style=prcpsty,numberlike=theorem]{principle}
\tcolorboxenvironment{principle}{colback=LightGreen}

\newcolumntype{L}[1]{>{\raggedleft\let\newline\\\arraybackslash\hspace{0pt}}m{#1}}
\newcolumntype{C}[1]{>{\centering\let\newline\\\arraybackslash\hspace{0pt}}m{#1}}
\newcolumntype{R}[1]{>{\raggedright\let\newline\\\arraybackslash\hspace{0pt}}m{#1}}

\hypersetup{
    pdftitle={Propuesta proyecto GTFS Transmilenio},
    pdfauthor={Alvarado Becerra Ludwig, Mahecha Reyes Albert, Ospina Rocha María Fernanda, Pinilla Naranjo Kevin Adrián},
    pdfsubject={Descripción del documento},
    pdfkeywords={palabras, interesantes},
    pdfcreator={overleaf2.ludwigalvarado.me},
    pdfproducer={pdflatex}
}


\setstretch{1.2}
\geometry{
    textheight=9in,
    textwidth=5.5in,
    top=1in,
    headheight=12pt,
    headsep=25pt,
    footskip=30pt
}

\lstdefinestyle{bashstyle}{
    language=bash,
    basicstyle=\ttfamily,
    backgroundcolor=\color{gray!10},
    keywordstyle=\color{blue},
    commentstyle=\color{green!40!black},
    stringstyle=\color{red},
    showstringspaces=false,
    numbers=left,
    numberstyle=\tiny\color{gray},
    breaklines=true,
    breakatwhitespace=true,
    frame=tb,
    rulecolor=\color{black!70},
    framerule=0.5pt,
    tabsize=4,
    captionpos=b
}

\lstdefinestyle{javastyle}{
    language=Java,
    basicstyle=\ttfamily,
    backgroundcolor=\color{gray!10},
    keywordstyle=\color{blue},
    commentstyle=\color{green!40!black},
    stringstyle=\color{red},
    showstringspaces=false,
    numbers=left,
    numberstyle=\tiny\color{gray},
    breaklines=true,
    breakatwhitespace=true,
    frame=tb,
    rulecolor=\color{black!70},
    framerule=0.5pt,
    tabsize=4,
    captionpos=b
}

\lstdefinestyle{pythonstyle}{
    language=Python,
    basicstyle=\ttfamily,
    backgroundcolor=\color{gray!10},
    keywordstyle=\color{blue},
    commentstyle=\color{green!40!black},
    stringstyle=\color{red},
    showstringspaces=false,
    numbers=left,
    numberstyle=\tiny\color{gray},
    breaklines=true,
    breakatwhitespace=true,
    frame=tb,
    rulecolor=\color{black!70},
    framerule=0.5pt,
    tabsize=4,
    captionpos=b
}

\lstdefinestyle{cppstyle}{
    language=C++,
    basicstyle=\ttfamily,
    backgroundcolor=\color{gray!10},
    keywordstyle=\color{blue},
    commentstyle=\color{green!40!black},
    stringstyle=\color{red},
    showstringspaces=false,
    numbers=left,
    numberstyle=\tiny\color{gray},
    breaklines=true,
    breakatwhitespace=true,
    frame=tb,
    rulecolor=\color{black!70},
    framerule=0.5pt,
    tabsize=4,
    captionpos=b
}



\pagestyle{fancy}
\fancyhf{}
\renewcommand{\headrulewidth}{0pt}
\fancyfoot[C]{\thepage}
\fancyfoot[C]{\footnotesize \thepage \quad \textbar{} \quad Este trabajo está bajo una licencia CC BY-SA 4.0.\\ Más info: \url{https://creativecommons.org/licenses/by/4.0/}}


% ------------------------------------------------------------------------------

\begin{document}

% ------------------------------------------------------------------------------
% Cover Page and ToC
% ------------------------------------------------------------------------------

\title{ \normalsize \textsc{}
	\\ [2.0cm]
	\HRule{1.5pt} \\
	\LARGE \textbf{\uppercase{FlowTM}
		\HRule{2.0pt} \\ [0.6cm] \LARGE{Agente racional para la optimización multiobjetivo de frecuencias y horarios de Transmilenio mediante GTFS} \vspace*{10\baselineskip}}
}
\date{}
\author{\textbf{Alvarado Becerra Ludwig} \\
\textbf{Mahecha Reyes Albert}\\
\textbf{Ospina Rocha María Fernanda} \\ 
\textbf{Pinilla Naranjo Kevin Adrián} \\ 
	 Inteligencia Artificial - 2026-2S}

\maketitle
\thispagestyle{empty}
\newpage

\tableofcontents
\thispagestyle{empty}
\newpage
\setcounter{page}{1}

\section{Descripción del problema}

TransMilenio, el sistema de transporte masivo de Bogotá, transporta diariamente cerca de dos millones de individuos por medio de su elemento troncal y zonal~\cite{transmilenio2024informe}. Sin embargo, la operación enfrenta un grave desbalance estructural entre los buses programados y la fluctuante demanda de pasajeros, sobre todo en las horas pico de la mañana y la tarde~\cite{transmilenio2024informe,bogotacomovamos2024encuesta}. El 49\% de los usuarios manifiesta un profundo desagrado ante las largas esperas en las plataformas, y el 49,2\% se queja de la congestión y los prolongados tiempos de viaje en los buses articulados, según la Encuesta Bogotá Cómo Vamos, ``Encuesta de Percepción Ciudadana 2024: Movilidad y Calidad de Vida en Bogotá,''~\cite{bogotacomovamos2024encuesta}

El sistema opera bajo una planificación de frecuencias y tablas horarias definida en su especificación estática GTFS (General Transit Feed Specification), la cual establece de forma fija la oferta de buses por ruta en cada franja horaria~\cite{transmilenio2026datosabiertos}. No obstante, la demanda real no se distribuye de manera uniforme entre las rutas que comparten un mismo corredor troncal, sino que varía de acuerdo con la hora, los patrones de origen-destino y la dinámica urbana~\cite{transmilenio2024informe,transmilenio2026datosabiertos}. 

Esta discrepancia entre una oferta fija y una demanda estocástica genera asimetrías operativas observables: mientras rutas de alta afluencia operan con sobreocupación crítica y tiempos de espera desbordados, rutas paralelas en el mismo corredor circulan con capacidad ociosa e intervalos subutilizados~\cite{bogotacomovamos2024encuesta,yang2020multiobjective}. Actualmente, la planificación del GTFS carece de mecanismos computacionales automatizados de ajuste que respondan con agilidad a estos patrones~\cite{yang2020multiobjective}. 

Ante la disponibilidad de datos abiertos oficiales publicados por TransMilenio (malla GTFS y registros históricos de validaciones horarias en estaciones)~\cite{transmilenio2026datosabiertos}, existe la oportunidad de abordar este problema mediante un Agente Racional. Este agente percibirá el estado de la red y la demanda, procesará la información mediante una función de utilidad multicriterio (balanceando tiempo de espera, nivel de ocupación y costo de flota) y actuará proponiendo ajustes a las frecuencias de despacho en el GTFS, respetando las restricciones operativas reales del sistema~\cite{yang2020multiobjective,russell2020artificial}. 



\section{Objetivos}

\subsection{General}

Diseñar e implementar un agente racional basado en inteligencia artificial que, a partir del análisis de datos históricos de oferta y demanda del sistema TransMilenio, proponga modificaciones a la planificación de frecuencias en la especificación GTFS con el fin de reducir el desequilibrio de ocupación entre rutas de un mismo corredor, sujeto a las restricciones operativas del sistema. 

 
\subsection{Específicos}

\begin{itemize}
    \item Formalizar el entorno y la estructura PEAS del agente mediante la estructuración del espacio de estados, acciones, supuestos y restricciones operativas a partir de los datos abiertos de la red GTFS y las validaciones horarias del Portal de Datos Abiertos de TransMilenio~\cite{transmilenio2026datosabiertos}. 
    \item Diseñar y calibrar una función de utilidad multicriterio que cuantifique matemáticamente el costo por tiempo de espera de los usuarios, el factor de sobreocupación vehicular y la penalización por uso ineficiente de la flota disponible.  
    \item Implementar un entorno de simulación de eventos discretos que represente el comportamiento de un corredor troncal piloto, permitiendo al agente interactuar con el entorno y evaluar el impacto de sus decisiones de despacho en horas pico. 
    \item Evaluar el desempeño del algoritmo de optimización del agente frente a la malla GTFS estática actual, verificando la convergencia del modelo y cuantificando la reducción porcentual en los tiempos de espera y el balance de carga en los buses. 
\end{itemize}


\section{Enlace vídeo}

Entrar con el correo \texttt{utadeo.edu.co}

\href{https://utadeoeduco0-my.sharepoint.com/:v:/g/personal/ludwig_alvaradob_utadeo_edu_co/IQCiziMBhBLFRr7USRWr229VAeKW1t1dJSIDF9qKpHm_7K4?nav=eyJyZWZlcnJhbEluZm8iOnsicmVmZXJyYWxBcHAiOiJPbmVEcml2ZUZvckJ1c2luZXNzIiwicmVmZXJyYWxBcHBQbGF0Zm9ybSI6IldlYiIsInJlZmVycmFsTW9kZSI6InZpZXciLCJyZWZlcnJhbFZpZXciOiJNeUZpbGVzTGlua0NvcHkifX0&e=Zf0X3v}{\textbf{Aquí}}


\addcontentsline{toc}{section}{Referencias}
\printbibliography

\newpage

\thispagestyle{empty}

\vspace*{0.3\textwidth}

\begin{figure}[H]
  \centering
  \includesvg[width=\textwidth]{img/LogoUtadeo.svg}
\end{figure}



\end{document}
